\documentclass[11pt,a4paper]{article}
\usepackage[margin=1.8cm]{geometry}
\usepackage[T1]{fontenc}
\usepackage[utf8]{inputenc}
\usepackage{listings,xcolor,hyperref,booktabs,enumitem,parskip}
\usepackage[expansion=false]{microtype}

\definecolor{leafgreen}{RGB}{27,116,76}
\definecolor{leafblue}{RGB}{30,94,135}
\definecolor{leafgray}{RGB}{245,247,246}
\hypersetup{colorlinks=true,linkcolor=leafgreen,urlcolor=leafblue,citecolor=leafgreen}
\lstset{
  basicstyle=\ttfamily\small,
  breaklines=true,
  backgroundcolor=\color{leafgray},
  frame=single,
  rulecolor=\color{black!20},
  xleftmargin=6pt,xrightmargin=6pt,
  aboveskip=7pt,belowskip=6pt,
  showstringspaces=false
}

\title{\textbf{LeafOS Resident Project Stack}\\[5pt]\large Quick Start \texttt{v0.5.0}}
\author{LeafOS}
\date{Updated 2026-07-21}

\begin{document}
\maketitle

\begin{center}
\fcolorbox{leafgreen}{leafgray}{\parbox{0.88\linewidth}{
\textbf{Current default:} attach the local stack to a project with
\texttt{leafctl live PATH}. The stack means all downloaded and locally
available models in this LeafOS instance. Models propose; CPU-side policy
executes and validates.
}}
\end{center}

\section*{1. Check The Host}

From PowerShell in the taskpack directory:

\begin{lstlisting}
.\bin\leafctl.ps1 doctor
.\bin\leafctl.ps1 versions
.\bin\leafctl.ps1 provider-stack check
\end{lstlisting}

The provider check validates the configured llama.cpp executable, local GGUF
model, Vulkan route, endpoint, and lifecycle settings. It does not download a
model.

\section*{2. Open A Project}

\begin{lstlisting}
.\bin\leafctl.ps1 live C:\R\MyProject
\end{lstlisting}

LeafOS performs bounded read-only intake, derives an objective when none is
supplied, creates or reattaches a persistent version-2 run, starts the resident
supervisor, and opens the terminal interface.

Inspect without starting work:

\begin{lstlisting}
.\bin\leafctl.ps1 live C:\R\MyProject --inspect --json
\end{lstlisting}

Create a minimal Catan2 project:

\begin{lstlisting}
.\bin\leafctl.ps1 live C:\Games\catan2 \
  --new --template catan2 --yes --budget 64
\end{lstlisting}

Intake does not add LeafOS control files to an existing target. Run state stays
in the LeafOS runs directory.

\section*{3. Give Instructions}

Inside the active window, plain language is enough:

\begin{lstlisting}
:improve
:improve add deterministic replay
:again
:make state changes easier to inspect
\end{lstlisting}

Every instruction converges on one lifecycle:

\begin{center}
\texttt{INTAKE -> PLAN -> APPROVE -> EXECUTE -> VALIDATE -> CHECKPOINT/REPORT}
\end{center}

The provider creates a bounded structured proposal. The CPU inlet validates
paths, commands, approval, dependencies, attempts, execution, and acceptance.
A failure preserves evidence and may create one bounded repair task.

\section*{4. Control The Resident Stack}

\begin{lstlisting}
:mode quiet
:mode auto
:targets cpu 80 gpu 90
:budget 64m
:pause
:resume
:drain
\end{lstlisting}

The utilization values are useful-work targets, not load quotas. Recent input,
responsiveness, RAM, VRAM, temperature, and provider health can reduce or close
dispatch. Pressing \texttt{q} closes only the interface. Drain finishes
accepted work. Stop uses a confirmed typed request.

\section*{5. Reconnect And Observe}

\begin{lstlisting}
.\bin\leafctl.ps1 tui --run active
.\bin\leafctl.ps1 loop status active
.\bin\leafctl.ps1 loop attach active --after 0
.\bin\leafctl.ps1 resident status active --json
\end{lstlisting}

Run directories retain the work order, queue, state, events, controls,
resident decision, provider artifacts, command output, validation, checkpoint,
telemetry, and report required to reconnect without conversation memory.

\section*{6. Verify}

\begin{lstlisting}
.\bin\leafctl.ps1 test visual --plain
.\bin\leafctl.ps1 test visual \
  --match provider --live-stack --plain
.\bin\leafctl.ps1 catan2-resident \
  --contract-only --iterations 3
\end{lstlisting}

The current discovery suite contains 203 tests. The live-stack demonstration
is explicitly real-provider evidence. Contract and fixture tests remain
clearly labeled and do not make that claim.

\section*{Authority And Privacy}

The TUI is a structured projection and typed control surface, not a shell or a
second scheduler. Models cannot approve themselves, write files directly, or
turn output into evidence. LeafOS displays plans, progress, decisions, tools,
validation, and accepted results. It does not expose or persist private model
chain-of-thought.

\section*{Next Documents}

\begin{tabular}{@{}p{0.31\linewidth}p{0.61\linewidth}@{}}
\toprule
Document & Purpose \\
\midrule
\texttt{DOCUMENTATION\_MAP.md} & Current authority and historical boundaries. \\
\texttt{ARCHITECTURE.md} & Process, data, provider, and UI boundaries. \\
\texttt{AGENTIC\_CLI.md} & Instruction normalization and lifecycle. \\
\texttt{RESIDENT\_STACK\_USAGE.md} & Resource policy, budgets, recovery, and soaks. \\
\texttt{MEDIUM\_MOE\_MODEL\_POLICY.md} & Model candidates, SSD truth, and qualification gates. \\
\texttt{TUI\_USAGE.md} & Pages, keys, snapshots, and control transport. \\
\bottomrule
\end{tabular}

\section*{Current Acceptance State}

The resident implementation, contract profile, 203-test observatory, native
TUI, bounded live-provider demonstration, and real four-minute Catan2 profile
pass. The real eight-minute and 64-minute soaks remain open evidence. A profile
is not complete merely because its command exists.

\section*{Optional Surfaces}

\begin{itemize}[nosep,leftmargin=1.4em]
  \item \texttt{oneshot} creates a handoff bundle; it does not run a project.
  \item \texttt{dashboard} displays the optional node/cluster subsystem.
  \item \texttt{node}, \texttt{nodes}, and \texttt{task} distribute explicit SSH jobs.
  \item Early agent-plan and graph commands remain compatibility paths.
\end{itemize}

Historical work orders and reports preserve their original terminology as
dated evidence. Use the Documentation Map when a historical command conflicts
with this guide.

\end{document}
